We implemented the FPPVS framework in Python and found that, as specified, it
could not be reproduced and performed poorly in our hands: about 10\%
acceptance, below random placement. The main causes were an unstated arrival
rate that overloaded the network, feasibility checks that came too late, a
sharing model that could overbook capacity, and an actor network that could not
tell data centers apart. We stated the workload as an explicit offered load,
added a look-ahead feasibility mask, an instance-overhead model for VNF sharing,
a per-data-center actor with a learnable placement prior, rollback-and-retry
admission and corrected federated training.

With these changes FPPVS accepts \fppvsMinAcc--\fppvsMaxAcc\% of requests on
both topologies at loads up to full capacity, above the 70\% target in every
main scenario and every seed. On the same scenarios, the original FPPVS as
specified in the paper reaches only \origMinAcc--\origMaxAcc\% and the 2022
adaptation \zhangMinAcc--\zhangMaxAcc\%, and we could not reproduce the paper's
published values with the original specification. Updated FPPVS also holds the
least bandwidth per running service of all five methods (about 57\% less than
greedy and random, 31\% less than the original FPPVS and 14\% less than the 2022
adaptation) and starts far fewer VNF instances than the methods without sharing.
The ablation study shows that VNF sharing and the problem
modifications drive these gains, while the learned part of the policy adds
little beyond its prior.

Future work: (1) evaluate on the real USAnet and Italian operator topologies and
traces once they are available; (2) re-implement the paper's baselines; (3) give
the agent a multi-request horizon, for example by treating a batch of arrivals
as one episode, so that learning can improve on the prior; and (4) add VNF
migration so that long-lived chains can be repacked when links fill up.
